\thispagestyle{empty}
\pagenumbering{arabic}
\renewcommand*{\thepage}{C-\arabic{page}}

\chapter{Technical documentation}

\section{Setup guide}

This guide was developed and tested on Arch Linux (kernel 6.18.24-1-lts) on an
11th Gen Intel Core i7-1165G7 @ 2.80GHz with 16 GiB RAM. For reference, the
benchmarks in this thesis were run on an Intel Xeon E312xx server with Ubuntu
24.04 and 64 GB RAM.

The workspace is pinned to \texttt{nightly-2025-08-19}. Stable Rust is not
enough, because the workspace uses nightly-only features. The guest programs are
compiled for \texttt{riscv32imac-unknown-none-elf}, while the prover runs on the
native Linux host target (\texttt{x86\_64-unknown-linux-gnu} here).

\begin{itemize}
	\item Rust toolchain \texttt{nightly-2025-08-19}.
	\item Guest target \texttt{riscv32imac-unknown-none-elf}.
	\item About 14 GiB of free RAM for the proving run.
\end{itemize}

Downloading \texttt{rustup} is sufficient, everything else can be downloaded through it.
After installing it, run:

\begin{minted}[breaklines,obeytabs=true,tabsize=2,breaksymbolleft=]{bash}
rustup toolchain install nightly-2025-08-19
rustup target add --toolchain nightly-2025-08-19 riscv32imac-unknown-none-elf
\end{minted}

\section{Proving pipeline}

The proving pipeline starts from a tiny guest ELF, executes it step by step in the VM, and folds each step into a CCS accumulator.

\subsection{Guest program}

The Fibonacci guest is intentionally small:

\begin{minted}[breaklines,obeytabs=true,tabsize=2,breaksymbolleft=]{rust}
#![no_main]
#![no_std]

guest::guest_main!(main);

fn main() {
    let n = 100;
    let mut a: u32 = 0;
    let mut b: u32 = 1;
    let mut sum: u32;
    for _ in 1..n {
        sum = a.wrapping_add(b);
        a = b;
        b = sum;
    }

    guest::write_result(b);
}
\end{minted}

It computes the 100th Fibonacci number with wrapping 32-bit arithmetic and writes the result through the guest runtime.

\subsection{Guest SDK}

The \texttt{guest} crate in \texttt{crates/guest/src/lib.rs} provides the runtime glue needed for \texttt{no\_std} RISC-V code.
This runtime exports the user \texttt{main}, sets up \texttt{\_start},
initializes \texttt{gp} and \texttt{sp}, installs a tiny allocator through an
\texttt{ecall}, and places the final answer at the fixed result address
consumed by the host.

\subsection{Compiled binary}

After compiling such a guest program with \texttt{cargo build --release --target riscv32imac-unknown-none-elf},
it can then be loaded into the ZkVM.

\subsection{VM execution and CCS}

The prover loads the ELF with \texttt{vm.load\_elf(...)} and executes it through \texttt{vm.run(...)}. Each callback receives the instruction trace together with the pre-state and post-state, register file, memory view, and raw code bytes. These are packed into \texttt{IVCStepInput}, arithmetized with \texttt{arithmetize(...)}, and then committed as the CCS witness for that step.

The CCS constraints are encoded manually in the Rust implementation by constructing sparse selector matrices, multisets, and coefficients for each instruction class. For the \texttt{JALR} case, the builder expands the relation into the following matrix row:

\begin{minted}[fontsize=\footnotesize,breaklines,obeytabs=true,tabsize=2,breaksymbolleft=]{rust}
/// Adds a JALR constraint that ensures the return address is written correctly:
/// z[IS_JALR] * (z[VAL_RD_OUT] - (z[PC_IN] + z[INSTRUCTION_SIZE])) = 0
fn jalr_constraint(&mut self) {
    let matrix_base_idx = self.matrices.len();

    // Matrix A: selects z[IS_JALR]
    let mut m_a = empty_sparse_matrix(self.m, self.layout.z_vector_size());
    m_a.coeffs[JALR_CONSTR].push((R::one(), self.layout.is_jalr_idx));

    // Matrix B: selects (z[VAL_RD_OUT] - z[PC_IN] - z[INSTRUCTION_SIZE])
    let mut m_b = empty_sparse_matrix(self.m, self.layout.z_vector_size());
    m_b.coeffs[JALR_CONSTR].push((R::one(), self.layout.val_rd_out_idx));
    m_b.coeffs[JALR_CONSTR].push((R::one().neg(), self.layout.pc_in_idx));
    m_b.coeffs[JALR_CONSTR].push((R::one().neg(), self.layout.instruction_size_idx));

    self.matrices.push(m_a);
    self.matrices.push(m_b);

    // Add multiset: A * B
    self.multisets
        .push(vec![matrix_base_idx, matrix_base_idx + 1]);
    self.coeffs.push(R::one());
}
\end{minted}

This enforces that the link register receives the correct return address. The same pattern repeats for constraining every instruction.

When the \texttt{debug} feature is enabled, the prover additionally runs \texttt{check\_relation\_debug(...)} on every step to ensure the witness satisfies the CCS relation, and \texttt{fold(...)} replays \texttt{NIFSVerifier::verify(...)} so the intermediate folding proof is checked before the next step is processed.

\section{How to run}

\begin{enumerate}
	\item Compile the Fibonacci guest.

	      \begin{minted}[breaklines,obeytabs=true,tabsize=2,breaksymbolleft=]{bash}
cd guests/fibonacci
cargo build --release --target riscv32imac-unknown-none-elf
\end{minted}

	      The resulting binary is stored at:

	      \begin{minted}[breaklines,obeytabs=true,tabsize=2,breaksymbolleft=]{text}
target/riscv32imac-unknown-none-elf/release/fibonacci
\end{minted}

	\item Run the zkVM prover from the repository root.

	      \begin{minted}[breaklines,obeytabs=true,tabsize=2,breaksymbolleft=]{bash}
RUST_LOG=debug cargo run --bin zkvm --features latticefold/parallel --release -- --guest ./target/riscv32imac-unknown-none-elf/release/fibonacci
\end{minted}
\end{enumerate}

The proving run is memory-intensive and should be done on a machine with roughly
14 GiB of free RAM or more.
